%==============================================================================
% Research Methods - Group assignment - pitch
% Groep 64 - Simon Borghart, Emiel De Pelsmaeker, Simon Boey
% Onderwerp binnen specialisatie: AI & Data Engineering
%==============================================================================

\documentclass[aspectratio=169]{beamer}

%==============================================================================
% Preamble
%==============================================================================

%---------- Layout ------------------------------------------------------------

\usetheme{hogent}
\usecolortheme{hgwhite}

%---------- Packages ----------------------------------------------------------

\usepackage[dutch]{babel}

\usepackage{booktabs}
\usepackage{multirow,multicol}
\usepackage{eurosym}

%---------- Metadata ----------------------------------------------------------

\title{Pitch onderzoeksvoorstel.}
\subtitle{AI \& Data Engineering -- Research Methods, 25-26, groep 64}
\author{Simon Borghart \and Emiel De Pelsmaeker \and Simon Boey}
\date{\today}

%==============================================================================
% Presentation content
%==============================================================================

\begin{document}

{
    \setbeamertemplate{background}[imgletter]
        {beverage-coffee-computer-984536.jpg}{P}

    \begin{frame}
        \maketitle
    \end{frame}
}

\begin{frame}
    \frametitle{Inhoud.}

    \tableofcontents
\end{frame}

%==============================================================================
% Deel 1: Probleem
%==============================================================================

\section{Probleem.}

\subsection{Context.}

{

\begin{frame}

    {\huge \textbf{Uber's volledige jaarbudget voor AI-tools was al op,
    nog voor het einde van het jaar.}}

    \bigskip

    \textbf{Praveen Neppalli Naga, CTO Uber, april 2026}

\end{frame}
}

\begin{frame}
    \frametitle{De trend: AI-agents in productie.}

    \begin{itemize}
        \item Bedrijven zetten steeds meer AI-agents in (bv.\ Claude Code,
            Cursor) voor echt productiewerk, niet enkel als chatbot.
        \item Elke stap die een agent zet kost tokens, en dus geld.
        \item Klassieke caching (context-, semantische caching) is
            ontworpen voor chatbots: die herhalen vaak gelijkaardige
            vragen.
        \item Een agent plant stap voor stap en hangt af van externe
            data en omgevingscontext, waardoor klassieke caching hier
            onvoldoende grip op heeft.
    \end{itemize}
\end{frame}

\begin{frame}
    \frametitle{Concreet voorbeeld: Uber.}

    \begin{itemize}
        \item Uber gaf in 2025 \textbf{3,4 miljard dollar} uit aan
            onderzoek en ontwikkeling.
        \item Toch raakte het volledige jaarbudget voor AI-tools op,
            nog voor het einde van het jaar (april 2026).
        \item Reden, volgens CTO Praveen Neppalli Naga: een sterke
            stijging van het gebruik van AI-coding-agents.
        \item Ongeveer \textbf{11\%} van de live backend-code bij Uber
            wordt intussen door AI-agents geschreven.
        \item Bron: Yahoo Finance/Benzinga, oorspronkelijk gerapporteerd
            door The Information.
    \end{itemize}
\end{frame}

\subsection{Onderzoek bevestigt het probleem.}

\begin{frame}
    \frametitle{Wat toont onderzoek aan?}

    \begin{itemize}
        \item Zhang et al.\ (2025, NeurIPS) tonen aan dat bestaande
            LLM-cachingtechnieken tekortschieten bij agenttoepassingen,
            precies om de reden hierboven: de output van een agent hangt
            af van externe data en omgevingscontext.
        \item Hierdoor blijven operationele kosten en latency onnodig
            hoog.
        \item Met hun eigen oplossing (\emph{Agentic Plan Caching})
            tonen zij een gemiddelde \textbf{kostenreductie van 50,31\%}
            en een \textbf{latencyreductie van 27,28\%} aan.
        \item Dit bevestigt zowel het probleem als de haalbaarheid van
            een oplossing.
    \end{itemize}
\end{frame}

\begin{frame}
    \frametitle{Probleemstelling.}

    Bedrijven die AI-agents inzetten in productieomgevingen kampen met
   \textbf{ hoge en moeilijk te voorspellen operationele kosten}. Klassieke
    cachingtechnieken houden onvoldoende rekening met het meerstaps- en
    contextafhankelijke karakter van agenttoepassingen. Daarnaast bestaan
    er aanvullende technieken zoals prompt-compressie en model-routing,
    maar een overzicht dat deze technieken specifiek afweegt voor
    agenttoepassingen ontbreekt. IT-professionals die agenttoepassingen
    ontwikkelen hebben daardoor weinig onderbouwde houvast bij het maken
    van een keuze.
\end{frame}

\subsection{Onderzoeksvraag.}

\begin{frame}
    \frametitle{Onderzoeksvraag.}

    \begin{quote}
        \textbf{Welke combinatie van kostenoptimalisatietechnieken
        (agentgeschikte caching, prompt-compressie en model-routing) is
        het meest geschikt voor een platformteam zoals dat van Uber, om
        de operationele kosten van AI-coding-agents in een
        productieomgeving te verlagen, zonder de kwaliteit van de output
        significant te verminderen?}
    \end{quote}
\end{frame}

\begin{frame}
    \frametitle{Doelgroep.}

    \begin{itemize}
        \item Platformteams en engineering-leads bij middelgrote tot
            grote softwareorganisaties.
        \item Specifiek: de teams die verantwoordelijk zijn voor de
            interne uitrol en het kostenbeheer van AI-coding-agents
            (bv.\ Claude Code, Cursor).
    \end{itemize}
\end{frame}

\begin{frame}
    \frametitle{Waarom dit onderwerp?}

    \small
    \begin{itemize}
        \item \textbf{Specific}: drie met naam benoemde technieken
            (agentgeschikte caching, prompt-compressie, model-routing)
            binnen \'e\'en concreet probleem.
        \item \textbf{Measurable}: kostenbesparing, latency en
            kwaliteitsbehoud zijn kwantificeerbaar, uit literatuur \'en
            uit een eigen PoC.
        \item \textbf{Achievable}: publiek beschikbare tools, frameworks
            en datasets, geen interne bedrijfsdata nodig.
        \item \textbf{Realistic/Relevant}: recent peer-reviewed onderzoek
            (NeurIPS 2025), verankerd in een publiek bedrijfsvoorbeeld
            (Uber).
        \item \textbf{Time-bound}: volledige aanpak uitgewerkt binnen
            twaalf weken.
    \end{itemize}
\end{frame}

%==============================================================================
% Deel 2: Aanpak
%==============================================================================

\section{Aanpak.}

\subsection{Deelvragen.}

\begin{frame}
    \frametitle{Deelvragen: probleemdomein.}

    \begin{itemize}
        \item Waarom schieten klassieke context- en semantische caching
            specifiek tekort bij LLM-agenttoepassingen?
        \item Welke factoren bepalen de kostenstructuur van
            LLM-agenttoepassingen (bv.\ aantal planningsstappen,
            promptlengte, modelkeuze per stap)?
        \item Welke risico's ontstaan wanneer kosten verlaagd worden
            zonder rekening te houden met de specifieke werking van
            agenttoepassingen?
    \end{itemize}
\end{frame}

\begin{frame}
    \frametitle{Deelvragen: oplossingsdomein.}

    \begin{itemize}
        \item Hoe werkt agentgeschikte caching (bv.\ het cachen van
            planstructuren in plaats van volledige antwoorden), en
            wanneer levert dit kostenbesparing op?
        \item In welke mate is prompt-compressie toepasbaar op de
            planningsstappen van een agent?
        \item Hoe werkt model-routing voor agenttoepassingen, en welke
            tools ondersteunen dit?
        \item Hoe kan de impact van deze technieken op kosten en
            kwaliteit gemeten en vergeleken worden bij
            agenttoepassingen?
    \end{itemize}
\end{frame}

\subsection{Methodologie.}

\begin{frame}
    \frametitle{Vier fasen, twaalf weken.}

    \small
    \begin{table}
        \begin{tabular}{@{}lcp{5.8cm}@{}}
        \toprule
        \textbf{Fase} & \textbf{Weken} & \textbf{Milestone} \\
        \midrule
        1. Probleemdomein  & 1--3   & Probleembeschrijving + kostenanalyse \\
        2. Literatuur       & 4--6   & Afwegingskader (caching, compressie, routing) \\
        3. Technisch luik   & 7--10  & PoC agentgeschikte caching: kost, latency, tokens \\
        4. Synthese         & 11--12 & Eindrapport \\
        \bottomrule
        \end{tabular}
    \end{table}
\end{frame}

\begin{frame}

    \begin{itemize}
        \item \textbf{Wat}: agentgeschikte caching (cachen van
            planstructuren) zelf implementeren en toepassen op een
            testscenario met een AI-coding-agent.
        \item \textbf{Meting}: tokenverbruik, kost per taak en latency,
            telkens met en zonder caching.
        \item \textbf{Vergelijking}: eigen resultaten afzetten tegen de
            kostenreductie (50,31\%) en latencyreductie (27,28\%) die
            Zhang et al.\ (2025) rapporteren.
        \item \textbf{Waarom dit telt}: vult het cesuurcriterium
            ``technische component'' in, naast het theoretische werk.
    \end{itemize}
\end{frame}

\subsection{Resultaat.}

\begin{frame}
    \frametitle{Verwacht resultaat.}

    \begin{itemize}
        \item Een afwegingskader waarmee een platformteam zoals dat van
            Uber kan kiezen tussen caching, prompt-compressie en
            model-routing.
        \item Concrete PoC-cijfers over kostenbesparing, latency en
            kwaliteitsbehoud, vergeleken met Zhang et al.\ (2025).
        \item Geen interne Uber-data nodig: publiek beschikbare tools,
            frameworks en datasets volstaan.
    \end{itemize}
\end{frame}

%==============================================================================
% Tot slot
%==============================================================================

\section{Tot slot.}

\begin{frame}
    \frametitle{Bronnen.}

    \small
    \begin{itemize}
        \item \textbf{Zhang et al.\ (2025)}, NeurIPS -- peer-reviewed
            conferentiebijdrage, kernprobleem en -oplossing (caching).
        \item \textbf{Chen et al.\ (2024)}, TMLR -- peer-reviewed
            tijdschriftartikel, \emph{FrugalGPT} (caching, routing,
            prompt-adaptatie).
        \item \textbf{Zhen et al.\ (2025)}, ACL/INLG -- peer-reviewed
            conferentiebijdrage, overzicht kostenefficiëntie.
        \item \textbf{Jiang et al.\ (2023)}, \textbf{Regmi \& Liu
            (2024--2025)}, \textbf{Dekoninck et al.\ (2024)} --
            arXiv-preprints over prompt-compressie, caching en
            model-routing.
        \item \textbf{Jain (2026)}, Yahoo Finance/Benzinga -- nieuws\-bericht,
            concreet bedrijfsvoorbeeld (Uber).
    \end{itemize}

    \medskip
    \begin{minipage}{.7\textwidth}
    Bewust een mix van brontypes: peer-reviewed publicaties,
    arXiv-preprints \'en een actueel nieuwsbericht.
    \end{minipage}
\end{frame}

\begin{frame}
    {\Huge \textbf{Vragen?}}

    \bigskip
    Bedankt voor jullie aandacht.
\end{frame}

\end{document}